\chapter{Conclusions and Recommendations}

This study developed and evaluated two CNN-based image representations for financial trend prediction, namely the baseline CNN--TIC model and the proposed CNN--MIC model. CNN--TIC used technical indicators only, while CNN--MIC extended the image representation by integrating macroeconomic indicators with technical indicators. Both models were evaluated under the same preprocessing pipeline, labeling strategy, validation procedures, and trading simulation rules across selected blue-chip stocks and penny stocks from the PSE market, as well as selected cryptocurrency assets.

The results showed that the inclusion of macroeconomic indicators can improve prediction performance in selected cases, but the advantage was not consistent across all assets and validation settings. CNN--MIC produced higher results in some stocks and cryptocurrencies, which suggests that macroeconomic indicators can contribute in terms of financial trend prediction. However, CNN--TIC remained more stable in many cases, such as when the technical structure of the asset appeared sufficient for prediction.The financial evaluation also showed that stronger classification results did not always lead to better trading returns. This indicates that predictive accuracy alone is not enough to produce stronger financial performance.

Overall, the study shows that adding macroeconomic indicators is a viable extension of the CNN--TIC framework, even if it was not consistent across all assets and evaluation settings. The main contribution of this work is the demonstration that macroeconomic indicators can be incorporated into an image-based CNN trading representation and may provide measurable benefits under specific market conditions.

Future work should address the dataset limitations by testing more assets, longer time frame for historical data, and adding more macroeconomic indicators that may contribute to its performance. Other image-based strategies and CNN architectures may also be explored to better capture the interaction between technical and macroeconomic signals. Lastly, future work may further analyze the model’s behavior when macroeconomic indicators are contributing most in improving prediction performance.